\documentclass[11pt]{article}
\usepackage{../../../tex/lecturestyle}

\title{Linear Recurrence Recovery (mod m)}
\author{}
\date{}

\begin{document}
\thispagestyle{fancy}

\lecturetitleband{Polynomials \& Formal Power Series}{Linear Recurrence Recovery (mod m)}{9}

\section{Overview}
Given the first $N$ terms of a sequence over a field, the Berlekamp--Massey
algorithm finds the shortest linear recurrence they satisfy: it scans the
terms left to right, and each time the current recurrence mispredicts a
term, it patches itself using the most recent earlier mismatch, dividing by
that mismatch's value. That division is the crux of why the algorithm needs
a field: it must be able to invert any nonzero element it encounters.

Competitive programming frequently works over $\mathbb{Z}/m\mathbb{Z}$ for a
modulus $m$ that is not prime --- $10^9$, a power of two, a product chosen
for a specific problem. When $m$ is composite, $\mathbb{Z}/m\mathbb{Z}$ has
zero divisors: nonzero elements with no inverse (e.g.\ $2$ modulo $4$).
Berlekamp--Massey's division step can hit exactly such an element and simply
gets stuck. The \textbf{Reeds--Sloane algorithm} (1985) solves the recovery
problem over $\mathbb{Z}/m\mathbb{Z}$ for arbitrary $m$, by working one prime
power at a time and, within a prime power, tracking several candidate
recurrences simultaneously instead of one.

This is the most involved lecture in the series; budget real time for it.

\section{Problem Statement}

\begin{definitionbox}
Fix a modulus $m \ge 1$ and a polynomial $S(x) = S_0 + S_1 x + \cdots +
S_{N-1}x^{N-1}$ over $\mathbb{Z}/m\mathbb{Z}$. A pair $(L, Q(x))$ is
\emph{admissible} if $Q(0) = 1$, $\deg Q \le L$, and
\[
  [x^i]\, Q(x) S(x) = 0 \qquad \text{for all } L \le i < N .
\]
The recovery problem asks for an admissible pair with $L$ minimal. Taking
$L = N$, $Q = 1$ is always admissible (the condition is vacuous), so a
solution always exists.
\end{definitionbox}

Writing $P(x) = Q(x)S(x) \bmod x^L$, admissibility is the same as asking that
$S(x) \equiv P(x)/Q(x) \pmod{x^N}$ as formal power series, where division by
$Q(x)$ makes sense because $Q(0) = 1$ is a unit --- $Q(x)$ has a formal power
series inverse over \emph{any} commutative ring with unit, not just a field.
This is the fact that keeps the problem meaningful even when the coefficient
ring is not a field.

\subsection{What breaks without a field}
Two things change when $m$ is composite.

First, as noted above, Berlekamp--Massey's patch step divides by the most
recent nonzero discrepancy, and a nonzero element of $\mathbb{Z}/m\mathbb{Z}$
need not be invertible.

Second, the uniqueness guarantee that makes the recovered recurrence
trustworthy --- over a field, if $2L \le N$ the admissible pair with minimal
$L$ is unique --- can fail outright over $\mathbb{Z}/m\mathbb{Z}$.

\begin{examplebox}[: non-uniqueness]
Take $m = 4$ and $S = (2,2,2,2,2,2,2,2)$. Both
\[
  Q_1(x) = 1 - x \qquad\text{and}\qquad Q_2(x) = 1 - 3x
\]
give an admissible pair with $L = 1$: both encode a first-order recurrence
($S_i = S_{i-1}$, respectively $S_i = 3S_{i-1}$), and both are satisfied,
since $3 \cdot 2 \equiv 2 \pmod 4$. So the minimal-$L$ solution is not
unique, even though $L$ is tiny compared to $N$.
\end{examplebox}

Fortunately, non-uniqueness of $Q(x)$ does not translate into ambiguity about
what the recurrence \emph{predicts}. The quantity that matters for
extrapolation is the power series $P(x)/Q(x)$, and that turns out to be
unique whenever $2L \le N$, by an argument that only needs $Q(0)$ to be a
unit --- not that the ring is a field.

\begin{theorembox}[(Uniqueness of the extrapolation, over any commutative ring)]
Let $R$ be a commutative ring with unit and $L$ a nonnegative integer with
$2L \le N$. If $(P_1, Q_1)$ and $(P_2, Q_2)$ both satisfy $\deg P_i \le L-1$,
$\deg Q_i \le L$, $Q_i(0) = 1$, and $S(x) \equiv P_i(x)/Q_i(x) \pmod{x^N}$,
then $P_1(x)/Q_1(x) = P_2(x)/Q_2(x)$ as formal power series.

\emph{Proof.} Clearing denominators in $P_1/Q_1 \equiv P_2/Q_2 \pmod {x^N}$
gives $P_1 Q_2 \equiv P_2 Q_1 \pmod{x^N}$. Both sides have degree at most
$2L - 1 < N$, so a congruence modulo $x^N$ between them is an exact equality
of polynomials: $P_1 Q_2 = P_2 Q_1$. Multiplying both sides by the formal
power series $(Q_1 Q_2)^{-1}$ (which exists since $Q_1(0) = Q_2(0) = 1$)
gives $P_1/Q_1 = P_2/Q_2$. \qed
\end{theorembox}

So even though the $m=4$ example has two different minimal recurrences, they
predict the same continuation (both say ``the next term is $2$''), and that
is all that is needed for the typical competitive-programming use of this
machinery: guess a recurrence from a prefix, then use it to jump to a huge
index $K$.

\section{Reducing to Prime Powers}
Factor $m = \prod_{j} p_j^{e_j}$. The plan is to solve the recovery problem
modulo each $p_j^{e_j}$ independently and recombine.

\begin{theorembox}[(CRT reduction)]
Let $(L, Q(x))$ be a minimal solution mod $m$, and let $(L_j, Q_j(x))$ be a
minimal solution mod $p_j^{e_j}$ for each $j$. Then $L = \max_j L_j$, and a
minimal solution mod $m$ is obtained by CRT-combining the $Q_j(x)$
coefficient-by-coefficient.

\emph{Proof sketch.} Reducing the mod-$m$ solution $(L,Q)$ modulo each
$p_j^{e_j}$ gives an admissible pair there (admissibility only involves
congruences, which survive reduction), so $L_j \le L$ for every $j$. Now CRT
the per-prime-power solutions together into a single $\tilde Q(x)$ of degree
$\max_j L_j$: this reduces correctly modulo each $p_j^{e_j}$, hence is
admissible modulo $m$ (admissibility mod a product of coprime moduli is
equivalent to admissibility mod each factor). So $L \le \max_j L_j$ too,
giving equality. \qed
\end{theorembox}

This turns the general problem into: \emph{solve the recovery problem modulo
a single prime power $p^e$.} That is where Reeds--Sloane's real content
lives.

\section{The Prime-Power Case}

Fix $m = p^e$. The obstruction identified earlier was division: Berlekamp--
Massey's patch step needs to divide by a discrepancy, and a discrepancy in
$\mathbb{Z}/p^e\mathbb{Z}$ might be a nonzero multiple of $p$, which has no
inverse.

The fix is to stop insisting on a single running candidate recurrence.
Instead, maintain $e$ of them in parallel, indexed by ``valuation level''
$j = 0, 1, \ldots, e-1$: slot $j$ holds the best recurrence found so far
among those with constant term \emph{exactly} $p^{j}$ (rather than $1$).
Every nonzero element of $\mathbb{Z}/p^e\mathbb{Z}$ decomposes uniquely as
$t \cdot p^{u}$ with $t$ a unit and $0 \le u < e$ (its $p$-adic valuation
$u$ and unit part $t$), so a discrepancy that cannot be inverted directly
can still be matched, exactly, against the constant term of \emph{some}
slot --- namely slot $k = e - 1 - u$. That match is what makes a valid
cancellation possible: dividing a unit by a unit is always fine.

\subsection{State}
For each slot $j = 0, \ldots, e-1$, after processing $n$ terms of $S$, we
maintain:
\begin{itemize}[leftmargin=1.5em]
  \item the current candidate $Q_j(x)$, with $[x^0]Q_j(x) = p^{j}$,
        $\deg Q_j \le L_j$, minimal $L_j$ such that $[x^i]Q_j(x)S(x) = 0$
        for $L_j \le i < n$;
  \item auxiliary data $(B_j(x), L_{B,j}, n_{B,j}, t_{B,j})$ recording the
        state of slot $j$ the last time $L_j$ increased --- this is exactly
        the ``most recent mismatch'' that ordinary Berlekamp--Massey keeps a
        single copy of, except here each slot keeps its own.
\end{itemize}
Initially $Q_j(x) = p^{j}$ (a constant), $L_j = 0$, and no auxiliary data
exists for any slot.

\subsection{One step}
At step $n$ (processing $S_n$), for every slot $j$ compute the discrepancy
$\Delta_j = [x^n]\,Q_j(x)S(x)$ under the \emph{current} $Q_j$. If
$\Delta_j = 0$, slot $j$ needs no change. Otherwise write
$\Delta_j = t_j\, p^{u_j}$ (unit times prime power), and let
$k = e - 1 - u_j$.

\begin{itemize}[leftmargin=1.5em]
\item If slot $k$ has no auxiliary data yet, there is nothing to cancel
      with: leave $Q_j(x)$ unchanged and set $L_j \leftarrow n+1$ (the
      recurrence is provisionally declared to need the whole prefix).
\item Otherwise, let $c = t_j \cdot t_{B,k}^{-1}$ --- valid since $t_{B,k}$
      is a unit by construction --- and update
      \[
        Q_j(x) \;\leftarrow\; Q_j(x) \;-\; c\, x^{\,n - n_{B,k}}\, B_k(x),
        \qquad
        L_j \;\leftarrow\; \max\bigl(L_j,\; L_{B,k} + n - n_{B,k}\bigr).
      \]
      This is structurally identical to the Berlekamp--Massey patch, with
      the donor $B_k$ drawn from a different slot's history rather than the
      same slot's own.
\end{itemize}

Finally, for every slot $j$ whose $L_j$ increased this step, snapshot the
\emph{old} state of slot $k = e-1-u_j$ (the donor used above, before this
step's updates) as the new auxiliary data for slot $j$:
\[
  B_j \leftarrow Q_k,\quad L_{B,j} \leftarrow L_k,\quad
  n_{B,j} \leftarrow n,\quad t_{B,j} \leftarrow t_k .
\]
Note the asymmetry: cancellation for slot $j$ borrows a \emph{donor}
indexed by $k = e-1-u_j$, but the auxiliary data saved \emph{into} slot $j$
records $k$'s own unit $t_k$, not $j$'s. This is what makes the recursion
close up on itself: an inductive argument (parallel to Berlekamp--Massey's
own minimality proof, done separately at each valuation level) shows the
pairing $j \leftrightarrow k$ is exactly self-consistent, so slot $j$'s
saved data always has the unit-multiple relationship needed for some later
slot to use it validly.

After all $N$ terms are processed, slot $0$ holds the answer: it is the one
whose constant term is $p^0 = 1$, matching the original problem.

\begin{notebox}
When $e = 1$ (i.e.\ $m$ is itself prime), there is only one slot, $j = 0$,
and $k = e - 1 - u_j$ is forced to equal $0$ as well: the algorithm
degenerates exactly to ordinary Berlekamp--Massey, matching against its own
history. The machinery above is genuinely only needed once $e \ge 2$.
\end{notebox}

\section{Complexity}
Given the factorization $m = \prod_j p_j^{e_j}$ with $E = \sum_j e_j$, a
direct implementation (recomputing each slot's discrepancy from scratch as a
length-$O(N)$ dot product) runs in
\[
  O\!\left(N^2 E + N E \log m\right)
\]
time: the $N^2 E$ term is the discrepancy computations across all slots and
all steps, and the $\log m$ factor accounts for the modular inverses.

\section{Reference Implementations}
Both examples below implement the full pipeline: factor $m$, run the
prime-power algorithm on each factor, and recombine with CRT.

\subsection{Rust}
\lstinputlisting[language=Rust]{example.rs}

\subsection{Python}
\lstinputlisting[language=Python]{example.py}

\section{Worked Examples}
\begin{examplebox}
Modulo $12 = 4 \times 3$, the Fibonacci-style sequence
$S_i = S_{i-1} + S_{i-2}$ (with $S_0 = S_1 = 1$) is recovered from ten terms
as $L = 2$, $Q(x) = 1 - x - x^2$ (i.e.\ coefficients $[1, 11, 11]$, since
$-1 \equiv 11 \pmod{12}$) --- exactly the expected characteristic
polynomial, obtained here purely from prefix data over a non-prime modulus.

For $S = (2,2,\ldots,2) \bmod 4$, the algorithm returns $Q(x) = 1 - x$ (one
of the two valid minimal recurrences from the non-uniqueness example above);
extrapolating with it continues to predict $2$ at every future position, as
guaranteed by the uniqueness-of-extrapolation theorem.
\end{examplebox}

\section{Further Reading}
\begin{itemize}[leftmargin=1.5em]
  \item AtCoder Algorithm Lectures index (source of this curriculum):
        \url{https://info.atcoder.jp/entry/algorithm_lectures/index}
  \item J. A. Reeds, N. J. A. Sloane, ``Shift-Register Synthesis (Modulo
        $m$)'', \emph{SIAM Journal on Computing} 14(3), 1985, pp.\ 505--513.
\end{itemize}

\end{document}
